%% §4 — the kernel made solid four ways, in lockstep with KERNEL.md §9,
%% at tag era6-campaign-4,
%% kernel/mechanization/README.md and the theorem names in
%% Kernel/Key.lean and Kernel/Trust.lean, kernel/tests/__init__.py,
%% and kernel/second/README.md. Line counts: wc -l at the tag.

\section{The kernel, made solid four ways}
\label{sec:kernel}

The kernel is the fixed part: the gate, the sealed runner, the result
order, the registry loader, and the driver that plays, regrades, and
renders --- five stdlib-only Python modules of 2,622 lines at the tag
(\code{driver.py} 689, \code{gate.py} 1,238, \code{registry.py} 191,
\code{results.py} 419, \code{runner.py} 85).
It is the one piece of code the gate cannot judge, because it is the
gate; it is generated like everything else; and so four disciplines
stand in for the gate (\cref{tab:disciplines}), each worded no
stronger than what it verifies. The kernel computes no truth of its
own: it runs judges, compares bytes, orders results, and draws the
log. A kernel bug can lose a result or misgrade one; that it can
never forge one --- no record reaches \grade{certified} unless an
interpreter run where the question lives passed --- is stated in the
specification and falsified by test (\code{test\_forge.py}): on a toy
registry bootstrapped from empty through the real gate, a search that
prints a witness together with the words \code{"grade": "certified"}
never becomes a result, because its witness does not replay; a bogus
certificate floors at \grade{claimed}; a broken carry-back only loses
the witness; and \code{regrade} lifts a stored proof to
\grade{certified} by check time alone.
\Cref{sec:faults} asks the same question of the real registry.

\begin{table}
\caption{The four disciplines that stand in for the gate on the
kernel, what each verifies, and what it does not. Sizes are lines at
the tag.}
\label{tab:disciplines}
\centering
\small
\begin{tabular}{@{}p{0.19\textwidth}p{0.44\textwidth}p{0.31\textwidth}@{}}
\toprule
discipline & verifies & does not verify \\
\midrule
proof --- Lean~4, 502 lines, core library only, toolchain pinned
  & the order on keys and the ratchet (\cref{thm:order}); the residual (\cref{thm:residual}); axiom audit printed at every build: order lemmas axiom-free, fold and residual lemmas on \textsf{propext} and \textsf{Quot.sound} through core list lemmas
  & the Python that computes keys and residuals; \code{test\_order.py} runs the same properties against \code{results.key} on random records \\
tests --- \code{python3 -m kernel.tests}, 2,259 lines, nine modules, a toy registry
  & the gate on the registry's own controls (every stamp re-derived, every shipped mutant refused again) and on the toy (every way of failing it fails); the seal (blank \code{PATH}, scratch directory, no tool reachable, a hit wall a result); pins recomputed from their definition; the frame contract; board and graph regenerated byte-identically; the fault harness
  & inputs no pinned run exercises beyond the toy; a stamp's agreement is a floor, since a pair may join after its language was stamped \\
second lineage --- \code{kernel/second/}, 727 lines, clean-room from the specification and the committed data
  & byte-agreement with the first kernel on \code{base}, \code{report}, and \code{graph} over every pinned run and a synthetic run holding the corners (\code{test\_second}); an audit hook on \code{open} confirms it reads no file of the first kernel
  & \code{play}, \code{admit}, \code{regrade}, which it never runs; eleven listed ambiguities of the specification on inputs no pinned run exercises \\
content pins --- sha256 over every entry, re-verified by the loader
  & that admitted bytes are the judged bytes; two entries claiming one name at one revision refused
  & anything about what the bytes mean \\
\bottomrule
\end{tabular}
\end{table}

\paragraph{The half that is mathematics is proved}
A \emph{key} is $(\textit{level}, \textit{bound}, \textit{grade},
\textit{gap})$ with $\textit{gap} \in \mathbb{N} \cup \{\bot\}$; write
$k \prec k'$ for the lexicographic order in which a smaller finite gap
is better and $\bot$ (no judgment ever ran) is worst, $k \preceq k'$
for $k \prec k'$ or $k = k'$, and $\textit{best}(\ell)$ for the fold of
a log $\ell$ that keeps the incumbent only while it is strictly
better; the best of an empty log is below every key, and $\ell\,\ell'$
is $\ell$ followed by $\ell'$.

\begin{theorem}[the order and the ratchet]
\label{thm:order}
$\prec$ is irreflexive, transitive, and asymmetric. For all logs
$\ell, \ell'$: $\textit{best}(\ell) \preceq
\textit{best}(\ell\,\ell')$; if $\textit{best}(\ell)$ is settled then
so is $\textit{best}(\ell\,\ell')$; and if the two have
the same level and bound, then the grade of the second is no lower
and, at equal grade, its gap no larger. A record that differs from
another only by a higher grade, or only by a smaller gap, is strictly
above it.
\end{theorem}

\begin{theorem}[the residual]
\label{thm:residual}
Let $h_1, \dots, h_n$ be the lineages of a route's hops, $j$ the
lineage of the judge that ran, and
$\textit{residual}(h, g, j) = \bigcup_{i \le g} h_i \cup j$. Then
$\textit{residual}(h, 0, j) = j$; $g \le g'$ implies
$\textit{residual}(h, g, j) \subseteq \textit{residual}(h, g', j)$;
for $g \le n$, appending hops upstream of the judgment leaves the
residual unchanged; and every element of a residual belongs to the
judge or to a hop, never to the search alone.
\end{theorem}

The first clause of \cref{thm:residual} is what makes
\grade{certified} route-independent; the last clause of
\cref{thm:order} is why a regrade, which keeps the value, is a strict
improvement exactly when it moves the grade or the gap.

\paragraph{The half that is operation is falsified, and doubled}
The test suite is the gate turned on the kernel; one of its tests is
how the blank \code{PATH} got into the specification --- an unset
\code{PATH} still found \code{/usr/bin/python3}. The second lineage is
the accelerator discipline of \cref{sec:generation} turned on the
kernel: two implementations of disjoint descent that agree on every
byte vouch for each other, as two results of disjoint lineage do. Its README records the
protocol --- what was read, what was run (the first kernel only as a
black box), one disclosure (an early traceback that named two source
files; the first kernel now refuses a corrupt registry without one)
--- and where they disagreed while being built (the key's tie rule,
the pin's exclusions, the binding rule for revisions) the
specification was made to say which is right. Manual review is the
last discipline, and the four above are what make it tractable: what
must be read is the trusted base, printed by \code{base}, and the
kernel's five modules, and the reviewer reads theorem statements,
tests, and a second implementation's disagreements rather than a
story.
